
\section{Automorfismos y morfismos de recubrimientos étales}
\label{V.3}

\begin{proposition}
\label{V.3.1}
Sea~$X$ étale, separado y de tipo finito sobre un preesquema~$Y$ localmente
noetheriano, y sea~$G$ un grupo finito que actúa sobre~$X$ por
$Y$-automorfismos. Entonces la acción de~$G$ es admisible y el preesquema
cociente~$X/G$ es étale sobre~$Y$.
\end{proposition}

No se supone que~$X$ sea finito sobre~$Y$; sin embargo, $X$ es
cuasiproyectivo sobre~$Y$, de donde se sigue la existencia de~$X/G$
por~\Ref{V.1.8}. Demostremos primero:

\begin{corollaire}
\label{V.3.2}
El morfismo $X\to X/G$ es étale.
\end{corollaire}

Puede suponerse, evidentemente, que~$G$ actúa transitivamente sobre el
conjunto de las componentes conexas de~$X$ y, después, considerando el
estabilizador de una componente conexa, que~$X$ es incluso conexo. Por último,
puede suponerse que la acción de~$G$ es fiel. Pero entonces se ve, como
en~\Ref{V.2.4}, que~$G$ actúa sin inercia; por~\Ref{V.2.3} se sigue que
$X\to X/G$ es étale. Se concluye gracias al lema siguiente.

\begin{lemme}[remordimiento por el Exposé I]
\label{V.3.3}
Sean $X\to X'\to Y$ morfismos de tipo finito, sea~$x$ un punto de~$X$ y sean
$x'$ e~$y$ sus imágenes. Supongamos que~$Y$ es localmente noetheriano. Si dos
de los tres morfismos considerados son étales en los puntos señalados, también
lo es el tercero.
\end{lemme}

Solo queda considerar el caso en que $X\to X'$ y $X\to Y$ son étales en~$x$,
y demostrar que $X'\to Y$ lo es en~$x'$ (que es el caso que necesitamos
\marginpar{117}
para~\Ref{V.3.1}). Mediante una extensión plana conveniente de la base~$Y$, se
reduce la cuestión al caso en que la extensión residual
\ifthenelse{\boolean{orig}}
{$\kappa(x)/\kappa(y)$}
{$\kres(x)/\kres(y)$}
es trivial. Consideremos los homomorfismos
$\cal{O}_y\to\cal{O}_{x'}\to\cal{O}_x$ y los homomorfismos obtenidos al pasar
a las completaciones. La hipótesis significa que
$\hat{\cal{O}_y}\to\hat{\cal{O}_x}$ y
$\hat{\cal{O}_{x'}}\to\hat{\cal{O}_x}$ son isomorfismos; por tanto,
$\hat{\cal{O}_y}\to\hat{\cal{O}_{x'}}$ también es un isomorfismo, lo que
demuestra el lema.

\begin{corollaire}
\label{V.3.4}
Si~$X$ es finito y étale sobre~$Y$, entonces~$X/G$ es finito y étale
sobre~$Y$.
\end{corollaire}

\begin{proposition}
\label{V.3.5}
Sean~$X$ y~$X'$ dos recubrimientos étales de~$Y$. Todo $Y$-morfismo
$f\colon X\to X'$ se factoriza como la composición de un morfismo étale
sobreyectivo $X\to X''$ y la inmersión canónica $X''\to X'$ de un subconjunto
$X''$ de~$X'$ que es a la vez abierto y cerrado.
\end{proposition}

Se sabe por~\Ref{I.4.8} que~$f$ es étale y, por tanto, un morfismo abierto. Por
otra parte, como~$X$ es finito sobre~$Y$, $f$ es cerrado; así,
$f(X)=X''$ es un subconjunto a la vez abierto y cerrado de~$X'$. Esto concluye
la demostración. (N.B.: habría bastado que~$X'$, en vez de ser un recubrimiento
étale, fuese no ramificado sobre~$Y$.)

\begin{corollaire}
\label{V.3.6}
Con las notaciones precedentes,
\ifthenelse{\boolean{orig}}
{$X\to X'$ es un epimorfismo estricto en la categoría de los preesquemas, y
$X'\to X''$ es un monomorfismo (incluso un monomorfismo estricto) en dicha
categoría}
{$X\to X''$ es un epimorfismo estricto en la categoría de los preesquemas, y
$X''\to X'$ es un monomorfismo (incluso un monomorfismo estricto) en dicha
categoría}.
\end{corollaire}

Por definición, la primera afirmación significa que la sucesión de morfismos
\begingroup
\shorthandoff{<>}% Spanish babel otherwise treats Xy-pic's bend delimiters as shorthands.
$$
\xymatrix@C=.7cm{X \times_{X''} X \ar@<2pt>[r]^-{\pr_1}
\ar@<-2pt>[r]_-{\pr_2} & X \ar[r] & X''}
$$
\endgroup
es exacta. Esto resulta de que $X\to X''$ es finito y fielmente plano, como se
ve fácilmente (\cf Grothendieck, \loccit). La afirmación dual para
$X''\to X'$ es aún más inmediata.

El Corolario~\Ref{V.3.6} nos será útil para la teoría del grupo fundamental en
el número siguiente. Quienes no gusten de la noción de epimorfismo estricto
pueden sustituir el Corolario~\Ref{V.3.6} por la variante que prefieran.
Aprovechemos solamente la ocasión para señalar que una factorización
\marginpar{118}
$f=f'f''$, con~$f''$ epimorfismo estricto y~$f'$ monomorfismo, es
necesariamente única salvo isomorfismo único (en cualquier categoría). Sin
embargo, puede existir al mismo tiempo una factorización $f=f_1f_2$ con las
propiedades duales: $f_2$ es un epimorfismo y~$f_1$ un monomorfismo estricto
(también única salvo isomorfismo único), que no sea isomorfa a la anterior.
Basta tomar, por ejemplo, la categoría de los espacios vectoriales topológicos
(separados, si así se desea) y un morfismo $u\colon X\to X'$ tal que~$u(X)$ no
sea cerrado.

\begin{proposition}
\label{V.3.7}
Sean~$Y$ un preesquema \emph{conexo} localmente noetheriano, $y$ un punto
de~$Y$ y~$\mathit{\Omega}$ una extensión algebraicamente cerrada
\ifthenelse{\boolean{orig}}
{de~$\kappa(y)$}
{de~$\kres(y)$}.
Para todo~$X$ sobre~$Y$, denotemos por~$X(\Omega)$ el conjunto de los puntos
de~$X$ con valores en~$\Omega$. Sean~$X$ y~$X'$ recubrimientos étales de~$Y$,
y sea $u\colon X\to X'$ un $Y$-morfismo tal que la aplicación correspondiente
$X(\Omega)\to X'(\Omega)$ sea una biyección. Entonces~$u$ es un isomorfismo.
\end{proposition}

Se reduce inmediatamente al caso en que~$X'$ es conexo. Como $X\to X'$ es
finito y étale, el número geométrico de puntos de una fibra de $X\to X'$ es
constante y es igual a~$1$ si y solo si el morfismo considerado es un
isomorfismo. Ahora bien, ese número es también el número de elementos de una
fibra de $X(\Omega)\to X'(\Omega)$, de donde se sigue la conclusión.
\ifthenelse{\boolean{orig}}{}
{\enlargethispage{.5cm}}%
